% Edited conversion of source p. 37.
\ReportHeading
  {Повзучість осесиметрично навантажених тіл обертання складної форми із функціонально-градієнтних матеріалів}
  {С.~М. Склепус}
  {Інститут механіки ім. С.~П. Тимошенка НАН України, Київ, Україна}
  {snsklepus@ukr.net}

Функціонально-градієнтні матеріали (ФГМ) утворюють клас композиційних матеріалів, у яких склад і структура змінюються за заданим законом, а разом із ними неперервно змінюються механічні та фізичні властивості. Через цю просторову неоднорідність класичні моделі для однорідних або кусково-однорідних матеріалів не можуть бути застосовані безпосередньо. Особливо недостатньо досліджено нелінійне деформування конструкцій із ФГМ за умов високотемпературної повзучості.

Розглянуто початково-крайову задачу визначення напружено-деформованого стану порожнистих осесиметрично навантажених тіл обертання складної форми, виготовлених із ФГМ. Методика просторового аналізу ґрунтується на спільному застосуванні методів R-функцій, Рітца та Рунге--Кутта--Мерсона. Крайову задачу зводять до варіаційної задачі для функціонала у формі Лагранжа; її розв’язують методом Рітца у поєднанні з методом R-функцій. Початкову задачу інтегрують методом Рунге--Кутта--Мерсона з автоматичним вибором кроку за часом.

Як числові приклади розраховано повзучість циліндра та тіла обертання зі складною формою меридіанного перерізу, виготовлених із ФГМ на основі алюмінію, армованого керамікою, під дією внутрішнього тиску. Об’ємна частка керамічних частинок змінювалася вздовж радіальної координати за лінійним і степеневим законами. Модуль Юнга визначався лінійною залежністю від об’ємної частки армувального матеріалу.

Розрахунки виконано для однорідного розподілу частинок, а також для ФГМ із меншим і більшим перепадом їх об’ємного вмісту вздовж радіуса. Ефективні властивості ізотропного ФГМ під час повзучості моделювали ізонапруженим підходом Рейсса та скалярною версією самоузгодженої моделі (\emph{Self-Consistent Model}, SCM). Показано, що модель Рейсса дає верхню оцінку деформацій повзучості, тобто систематично завищує їх значення.
